\documentclass[10pt,letterpaper]{article}

% --- Configuración general ---
\usepackage[letterpaper,margin=1.35cm]{geometry}
\usepackage[T1]{fontenc}
\usepackage[utf8]{inputenc}
\usepackage[spanish,es-nodecimaldot]{babel}
\usepackage{lmodern}
\usepackage{enumitem}
\usepackage{tabularx}
\usepackage{array}
\usepackage{xcolor}
\usepackage[hidelinks]{hyperref}
\usepackage{titlesec}
\usepackage{microtype}

\pagestyle{empty}
\setlength{\parindent}{0pt}
\setlength{\parskip}{0pt}
\setlength{\tabcolsep}{0pt}
\renewcommand{\familydefault}{\sfdefault}

% Ajustes de secciones
\titleformat{\section}
  {\large\bfseries}
  {}{0pt}{\MakeUppercase}[\vspace{-3pt}\titlerule]

\titlespacing*{\section}{0pt}{8pt}{5pt}

% Viñetas compactas
\setlist[itemize]{
  leftmargin=1.25em,
  itemsep=1pt,
  topsep=2pt,
  parsep=0pt,
  partopsep=0pt
}

\newcommand{\cvheading}[4]{%
  \begin{tabularx}{\textwidth}{@{}X>{\raggedleft\arraybackslash}p{4.0cm}@{}}
    \textbf{#1} & \textbf{#2}\\
    \textit{#3} & \textit{#4}
  \end{tabularx}
}

\begin{document}

% --- Encabezado ---
\begin{center}
  {\LARGE\textbf{Miguel Angel Montoya Cruz}}\\[4pt]
  Ciudad Valles, San Luis Potosí, México
  \enspace$\cdot$\enspace
  \href{https://linkedin.com/in/mangelmcz}{linkedin.com/in/mangelmcz}
  \href{mailto:miguelmontoya1015@icloud.com}{miguelmontoya1015@icloud.com}
  \enspace$\cdot$\enspace
  \href{https://github.com/}{GitHub}
  \enspace$\cdot$\enspace
  \href{https://orcid.org/}{ORCID}
\end{center}

\vspace{2pt}
Estudiante de último semestre de Ingeniería en Sistemas Computacionales con interés en desarrollo de software y programación. Cuento con experiencia en proyectos académicos y personales utilizando diferentes tecnologías. Busco una oportunidad de residencia profesional para aplicar mis conocimientos, adquirir experiencia y contribuir al desarrollo de soluciones tecnológicas.

\section*{Proyectos}

\cvheading
  {Aplicación móvil de finanzas con modelo de Inteligencia Artificial}
  {Ciudad Valles, San Luis Potosí}
  {Proyecto académico/personal}
  {Feb. 2026 -- May. 2026}
\begin{itemize}
  \item Desarrollo del frontend móvil con Flutter y del backend con FastAPI.
  \item Entrenamiento de modelos de regresión y clasificación usando Scikit-learn.
  \item Manipulación y procesamiento de datos mediante las librerías de Python NumPy y Pandas.
  \item Desarrollo de visualizaciones de datos mediante la librería de Python Matplotlib.
  \item Implementación de la metodología ágil Kanban para la organización y seguimiento del proyecto.
\end{itemize}

\vspace{3pt}
\cvheading
  {Biblioteca Digital Colaborativa}
  {Ciudad Valles, San Luis Potosí}
  {Proyecto académico}
  {Sep. 2025 -- Dic. 2025}
\begin{itemize}
  \item Desarrollo de una aplicación web utilizando React, Node.js y PostgreSQL.
  \item Implementación de Docker para la ejecución y el despliegue de la aplicación mediante contenedores.
  \item Desarrollo de funcionalidades de autenticación de usuarios, operaciones CRUD, votaciones, comentarios, filtros avanzados y panel de usuario.
  \item Implementación de la metodología ágil Kanban para la gestión del proyecto.
  \item Colaboración en la elaboración de un artículo académico sobre el proyecto.
\end{itemize}

\vspace{3pt}
\cvheading
  {Personal English School}
  {Ciudad Valles, San Luis Potosí}
  {Programador web}
  {Oct. 2024}
\begin{itemize}
  \item Desarrollo y diseño de la página web de la empresa utilizando HTML, CSS y JavaScript.
\end{itemize}

\section*{Educación}

\cvheading
  {TECNM / Instituto Tecnológico de Ciudad Valles}
  {Ciudad Valles, San Luis Potosí}
  {Ingeniería en Sistemas Computacionales}
  {Agosto 2022 -- Presente}

\section*{Habilidades técnicas}
\begin{itemize}
  \item \textbf{Idiomas adicionales:} Inglés B1.
  \item \textbf{Lenguajes:} Python, C, C++, HTML, CSS y JavaScript.
  \item \textbf{Frameworks:} Flutter y React.
  \item \textbf{Backend:} Node.js, FastAPI y Flask.
  \item \textbf{Bases de datos:} SQLite y PostgreSQL.
  \item \textbf{Herramientas:} Git, GitHub, Docker, Postman, LaTeX, SimulIDE y GNU Octave (software matemático).
  \item \textbf{Sistemas operativos:} Windows, macOS y Linux.
  \item \textbf{Hardware:} Programación y creación de proyectos electrónicos con Arduino Uno y microcontroladores PIC en lenguaje ensamblador.
\end{itemize}

\end{document}
